\section{Visao Geral}

O Graph500 e um benchmark voltado para avaliar sistemas de alto desempenho em
cargas de trabalho centradas em grafos. Diferentemente de benchmarks tradicionais
como o LINPACK, que enfatizam operacoes numericas densas e desempenho em ponto
flutuante, o Graph500 foi definido para representar aplicacoes intensivas em
dados, nas quais o desempenho depende de acesso irregular a memoria,
comunicacao, construcao de estruturas esparsas e travessia de grafos
\cite{murphygraph500, graph500spec}.

A especificacao do Graph500 define um conjunto de problemas sobre um grafo
sintetico nao direcionado e ponderado. O grafo e inicialmente gerado como uma
lista de tuplas de arestas. A partir dessa lista, o benchmark mede a construcao
de uma estrutura interna de grafo e, em seguida, executa kernels de analise sobre
essa mesma estrutura. A versao da especificacao usada neste trabalho descreve os
benchmarks de busca (\textit{Search}) e caminhos minimos de fonte unica
(\textit{Shortest Path}) \cite{graph500spec}.

O objetivo do benchmark nao e apenas medir o tempo de uma implementacao isolada
de BFS ou SSSP. Ele procura exercitar o sistema em um fluxo mais proximo de
aplicacoes reais de grafos: gerar uma entrada irregular, construir uma estrutura
de dados adequada, executar varias consultas a partir de vertices diferentes,
validar os resultados e reportar estatisticas de desempenho. Por isso, o
Graph500 e particularmente relevante para este trabalho, que estuda travessia de
grafos em GPU e comunicacao entre unidades de processamento.

\section{Fluxo Geral do Benchmark}

A execucao do Graph500 e organizada em etapas. Primeiro, o benchmark gera uma
lista de arestas. Em seguida, o Kernel 1 constroi uma representacao de grafo a
partir dessa lista. Depois sao escolhidos ate 64 vertices de origem, chamados de
\textit{search keys}, que serao usados nas execucoes de BFS e SSSP. Para cada
origem, o Kernel 2 executa uma BFS e produz uma arvore de predecessores. Em outro
laco, o Kernel 3 executa SSSP e produz predecessores e distancias. Por fim, os
resultados sao validados e as metricas de desempenho sao calculadas
\cite{graph500spec}.

Nem todas as etapas entram no tempo reportado. A geracao da lista de arestas e a
selecao das raizes nao sao cronometradas. O tempo medido inclui a construcao do
grafo no Kernel 1 e as execucoes dos kernels de travessia. A validacao tambem nao
e incluida no tempo de desempenho, mas e obrigatoria para demonstrar que os
resultados produzidos pelas implementacoes sao corretos.

Uma caracteristica importante da especificacao e que os kernels 2 e 3 sao
executados em lacos separados. A execucao de SSSP nao pode reutilizar informacoes
computadas pela BFS, e diferentes invocacoes de um mesmo kernel tambem nao podem
compartilhar informacoes entre si. Essa regra evita que uma implementacao se
beneficie de conhecimento acumulado em buscas anteriores e torna a medicao mais
representativa da execucao de consultas independentes.

\section{Geracao do Grafo}

O Graph500 utiliza um gerador sintetico escalavel baseado em grafos de Kronecker,
semelhante ao modelo R-MAT. A entrada inicial do benchmark e uma lista de tuplas
de arestas. Cada tupla contem dois vertices, \textit{StartVertex} e
\textit{EndVertex}, e opcionalmente um peso. Quando apenas o Kernel 2, de BFS, e
executado, a especificacao permite nao gerar pesos, pois eles nao sao usados pela
busca em largura e aumentariam o consumo de memoria sem necessidade
\cite{graph500spec}.

O benchmark recebe como parametro principal o valor \texttt{SCALE}. O numero de
vertices do grafo e definido como:

\begin{equation}
N = 2^{\mathrm{SCALE}}
\end{equation}

O numero de arestas e definido a partir do \texttt{edgefactor}, cujo valor padrao
na especificacao e 16:

\begin{equation}
M = \mathrm{edgefactor} \times N
\end{equation}

Assim, aumentar \texttt{SCALE} em uma unidade dobra o numero de vertices e,
mantido o mesmo \texttt{edgefactor}, tambem dobra o numero de arestas geradas. A
Tabela~\ref{tab:graph500-classes} mostra as classes de problema definidas na
especificacao.

\begin{table}[H]
\centering
\begin{tabular}{l|c|c|c}
Classe & SCALE & Edge factor & Tamanho aproximado para BFS com 64 bits/aresta \\
\hline
Toy & 26 & 16 & 0,017 TB \\
Mini & 29 & 16 & 0,137 TB \\
Small & 32 & 16 & 1,100 TB \\
Medium & 36 & 16 & 17,592 TB \\
Large & 39 & 16 & 140,738 TB \\
Huge & 42 & 16 & 1125,900 TB \\
\end{tabular}
\caption{Classes de problema do Graph500 segundo a especificacao.}
\label{tab:graph500-classes}
\end{table}

O gerador R-MAT pode ser descrito como uma divisao recursiva da matriz de
adjacencia. Para cada aresta, a matriz e dividida em quatro quadrantes, e o
gerador escolhe em qual quadrante a aresta sera posicionada usando probabilidades
desiguais. A especificacao define os parametros iniciadores como
$A = 0{,}57$, $B = 0{,}19$, $C = 0{,}19$ e
$D = 1 - (A+B+C) = 0{,}05$ \cite{graph500spec}. Esse processo e repetido ate que
os bits dos identificadores dos vertices da aresta sejam definidos.

Essa distribuicao gera grafos com estrutura irregular, contendo vertices de grau
muito alto e muitos vertices de grau baixo. Esse comportamento e desejavel para o
benchmark, pois se aproxima de propriedades observadas em grafos reais e cria
desafios para balanceamento de carga, acesso a memoria e comunicacao. O gerador
tambem pode produzir auto-lacos e arestas multiplas entre os mesmos vertices.
Esses elementos devem fazer parte da lista de entrada fornecida ao Kernel 1,
mesmo que possam ser ignorados ou tratados posteriormente pelas estruturas usadas
nos kernels de travessia.

Depois da geracao, a especificacao exige que os identificadores dos vertices
sejam permutados aleatoriamente e que a lista de arestas tambem seja embaralhada.
Esse passo e importante porque remove localidade artificial introduzida pelo
processo de geracao. Sem essa permutacao, uma implementacao poderia explorar
padroes de localidade que nao estariam necessariamente presentes em grafos reais.
Em implementacoes paralelas, a numeracao dos vertices deve permanecer globalmente
consistente, e a distribuicao dos dados nao deve introduzir localidade indevida
entre processadores \cite{graph500spec}.

\section{Kernel 1: Construcao do Grafo}

O Kernel 1 transforma a lista de tuplas de arestas em uma estrutura de grafo usada
pelos kernels seguintes. A especificacao nao exige uma representacao especifica.
Uma implementacao pode usar matrizes esparsas, listas de adjacencia, CSR
(\textit{Compressed Sparse Row}), estruturas distribuidas ou outras
representacoes adequadas ao hardware alvo. A unica exigencia e que a estrutura
construida seja suficiente para executar os kernels posteriores e validar seus
resultados \cite{graph500spec}.

Esse kernel recebe apenas a lista de arestas e o tamanho dessa lista. Informacoes
como o numero total de vertices devem ser determinadas pela propria
implementacao, se forem necessarias. O grafo resultante deve representar uma
entrada nao direcionada. Portanto, uma aresta entre $u$ e $v$ deve permitir a
travessia em ambas as direcoes. Auto-lacos podem ser removidos na estrutura
interna, mas a implementacao precisa respeitar as regras de contagem e validacao
do benchmark.

O tempo de construcao do grafo e medido. Isso e relevante porque, em aplicacoes
reais, o custo de transformar uma lista de relacoes em uma estrutura eficiente
para consulta pode ser significativo. Em ambientes paralelos, essa etapa tambem
inclui custos como particionamento, redistribuicao de arestas, calculo de graus,
alocacao de memoria e montagem de listas de adjacencia.

Depois de construido, o grafo nao pode ser modificado entre os kernels de forma a
beneficiar um kernel especifico. A especificacao permite reservar espaco para
marcacoes temporarias ou travas, desde que esse espaco seja usado apenas pela
execucao do kernel atual. Essa regra separa a estrutura permanente do grafo dos
dados temporarios de cada travessia.

\section{Selecao das Raizes}

Antes das travessias, o Graph500 seleciona aleatoriamente ate 64 vertices de
origem. Esses vertices devem ter grau pelo menos um, desconsiderando auto-lacos.
A motivacao e evitar buscas triviais iniciadas em vertices isolados, que nao
exercitariam o sistema nem representariam adequadamente o custo de percorrer o
grafo.

Se o grafo gerado possuir menos de 64 vertices nao isolados, a especificacao
permite executar menos buscas. Em grafos nos tamanhos normalmente usados pelo
benchmark, esse caso e improvavel, mas a regra existe porque o gerador e
probabilistico. A selecao das raizes nao entra no tempo medido, mas o numero de
buscas efetivamente executado deve ser informado na saida.

\section{Kernel 2: Busca em Largura}

O Kernel 2 executa uma Busca em Largura, ou BFS, a partir de cada raiz escolhida.
A BFS visita o grafo por niveis: primeiro a origem, depois seus vizinhos, depois
os vizinhos ainda nao visitados desses vertices, e assim por diante. Ao final, o
kernel deve produzir uma arvore de predecessores. O pai da raiz deve ser a
propria raiz, e vertices nao alcancados devem receber predecessor igual a -1
\cite{graph500spec, cormenalgorithms}.

A especificacao nao fixa o algoritmo interno de BFS. Uma implementacao pode usar
fila, fronteiras, bitmaps, algoritmos direcionais, particionamento distribuido ou
estrategias especificas para GPU, desde que o resultado final seja uma arvore de
BFS valida. Essa liberdade e importante porque arquiteturas diferentes exigem
estrategias diferentes para obter desempenho.

O padrao de acesso a memoria do Kernel 2 e uma das principais razoes para a
existencia do Graph500. A BFS em grafos esparsos possui acesso dependente dos
dados, baixa previsibilidade e pouca computacao por aresta. Alem disso, grafos
gerados por R-MAT possuem vertices de grau muito variavel, criando
desbalanceamento de carga. Em sistemas paralelos, varias threads ou processos
podem tentar descobrir o mesmo vertice ao mesmo tempo, gerando contencao e
necessidade de sincronizacao.

A especificacao permite uma condicao de corrida benigna na escolha do pai de um
vertice. Quando vertices do nivel $k$ descobrem vertices do nivel $k+1$, nao e
necessario garantir que o primeiro visitante seja sempre registrado como pai. Se
mais de um predecessor valido poderia descobrir o mesmo vertice no nivel correto,
qualquer um deles pode ser aceito. O requisito e que a arvore final seja uma
arvore de BFS correta.

\section{Kernel 3: Caminhos Minimos de Fonte Unica}

O Kernel 3 executa SSSP (\textit{Single Source Shortest Paths}), isto e, calcula
a menor distancia da raiz para todos os vertices alcancaveis em um grafo com
pesos nao negativos. A saida inclui tanto um vetor de distancias quanto uma
arvore de predecessores. Assim como na BFS, o pai da origem deve ser a propria
origem, e vertices nao alcancados devem ter predecessor igual a -1
\cite{graph500spec, cormenalgorithms}.

Esse kernel e separado da BFS e nao pode reutilizar dados produzidos pelo Kernel
2. Embora muitos algoritmos de SSSP tenham estrutura parecida com uma travessia
por fronteiras, o uso de pesos adiciona mais trabalho e mais dependencias. A
especificacao tambem permite corridas benignas em SSSP, desde que as distancias e
a arvore final representem caminhos minimos corretos.

Neste trabalho, a enfase principal esta na BFS, pois ela corresponde ao Kernel 2
do Graph500 e e o nucleo da implementacao estudada. Ainda assim, a presenca do
Kernel 3 na especificacao e importante porque mostra que o benchmark foi pensado
como uma familia de analises sobre a mesma estrutura de grafo, nao apenas como
uma unica rotina de busca.

\section{Validacao}

Em escalas grandes, nao e pratico validar o resultado comparando-o com uma saida
padrao fixa. O grafo e gerado pseudoaleatoriamente, e diferentes implementacoes
paralelas podem escolher predecessores diferentes, mas igualmente validos. Por
isso, o Graph500 usa uma validacao baseada em propriedades do resultado
\cite{graph500spec}.

Para BFS, a validacao verifica se a arvore de predecessores realmente forma uma
arvore sem ciclos, se a raiz aponta para si mesma, se vertices nao visitados estao
fora da componente alcancavel da raiz, se cada vertice visitado possui uma aresta
real ligando-o ao seu pai e se os niveis da arvore respeitam a propriedade da
BFS. Em particular, para qualquer aresta do grafo, os niveis dos seus extremos
nao podem diferir por mais de um quando ambos pertencem a arvore.

Para SSSP, a validacao tambem verifica a consistencia da arvore de predecessores,
mas substitui a propriedade de niveis por propriedades de distancia. As
distancias associadas aos vertices devem respeitar os pesos das arestas, e a
arvore deve representar caminhos minimos a partir da origem. A validacao e
executada apos cada busca, mas seu tempo nao entra nas metricas de desempenho.

Essa separacao entre tempo de execucao e validacao e essencial para o benchmark.
Ela permite que uma implementacao seja avaliada por desempenho, mas impede que
resultados incorretos sejam reportados como validos. Em implementacoes paralelas,
essa validacao tambem e importante porque erros de sincronizacao podem produzir
arvores aparentemente plausiveis, mas que violam propriedades da BFS ou do SSSP.

\section{Metricas de Desempenho}

A principal metrica do Graph500 e TEPS (\textit{Traversed Edges Per Second}), ou
arestas atravessadas por segundo. Essa metrica foi escolhida porque a BFS e o
SSSP em grafos esparsos executam pouca aritmetica por aresta. Portanto, medir
operacoes de ponto flutuante por segundo nao representa adequadamente o gargalo
do problema. TEPS mede a taxa com que o sistema percorre arestas e, assim,
reflete melhor largura de banda efetiva, latencia de memoria, comunicacao,
sincronizacao e balanceamento de carga \cite{graph500spec}.

Para uma execucao de um kernel, a taxa e definida como:

\begin{equation}
\mathrm{TEPS} = \frac{m}{t}
\end{equation}

onde $m$ e o numero de arestas atravessadas na componente visitada e $t$ e o
tempo medido da execucao. Para arestas nao correspondentes a auto-lacos, a
contagem considera o numero de tuplas da componente dividido por dois, pois o
grafo e nao direcionado. Auto-lacos sao contabilizados conforme as regras da
especificacao.

Como cada benchmark executa varias buscas, a saida deve incluir estatisticas dos
tempos, das arestas percorridas e das taxas TEPS. A especificacao solicita
valores como minimo, primeiro quartil, mediana, terceiro quartil, maximo, media e
desvio padrao. Para TEPS, por ser uma taxa, a media reportada deve ser harmonica,
nao aritmetica. Essa escolha evita distorcoes ao combinar taxas obtidas em buscas
com tempos diferentes.

Os campos de saida incluem, entre outros, \texttt{SCALE}, \texttt{edgefactor},
\texttt{NBFS}, \texttt{construction\_time}, estatisticas de tempo da BFS,
estatisticas de quantidade de arestas atravessadas e estatisticas de TEPS. Quando
apenas um dos kernels e executado, os campos de TEPS do outro kernel podem ser
zerados, conforme permitido pela especificacao \cite{graph500spec}.

\section{Criterios de Avaliacao}

A especificacao define criterios de avaliacao em ordem aproximada de importancia.
O primeiro criterio e aderir de forma justa ao objetivo do benchmark. Em seguida,
considera-se o maior tamanho de problema que uma maquina consegue executar e o
menor tempo de execucao para um tamanho de problema dado. Criterios como tamanho
do codigo, tempo de desenvolvimento, manutenibilidade e extensibilidade tambem
sao desejaveis, mas aparecem como objetivos secundarios \cite{graph500spec}.

Essa ordem de criterios reforca a natureza do Graph500. O benchmark nao procura
premiar apenas implementacoes pequenas ou simples, mas sim implementacoes que
respeitem a especificacao, processem grafos grandes e executem travessias de
forma eficiente. Para este trabalho, isso significa que desempenho em GPU so e
relevante quando acompanhado de construcao correta do grafo, selecao adequada de
raizes, preservacao da semantica da BFS e validacao dos predecessores gerados.


